\hypertarget{chapter-9-modern-font-technology-with-lualatex-and-xelatex}{%
\chapter{Modern Font Technology with LuaLaTeX and XeLaTeX}\label{chapter-9-modern-font-technology-with-lualatex-and-xelatex}}

For most of LaTeX's history, using a font meant using one of a small set of Type 1 PostScript fonts specially prepared for TeX, each requiring its own metric files and encoding support before it could be used at all. This is still how the original pdfTeX engine works, and it is why pdfTeX-based documents are limited to a comparatively small library of fonts unless someone has already done the preparation work. LuaLaTeX and XeLaTeX exist in part to remove this limitation: both can load ordinary OpenType and TrueType font files directly, the same font files used by any other modern application, with no separate preparation step required. This chapter covers what that capability actually makes possible and where the two engines still differ enough to matter.

\hypertarget{why-not-pdflatex}{%
\section{Why Not pdfLaTeX}\label{why-not-pdflatex}}

pdfTeX remains the fastest of the three engines and is still the right choice for documents that only need Computer Modern or one of the small set of professionally prepared Type 1 fonts already packaged for TeX. For anything else, and especially for a book-length project intended to use a specific, deliberately chosen typeface, pdfTeX is the wrong default. Using an arbitrary OpenType font under pdfTeX requires generating Type 1 metric files for it first, a process that is at best an extra build step and at worst does not work cleanly for a given font at all, particularly for fonts with modern OpenType features that have no equivalent in the older metric format pdfTeX expects.

LuaLaTeX and XeLaTeX both sidestep this entirely: a font is specified by name or by file path, loaded directly by the engine, and its full feature set becomes available through \pkg{fontspec} without any separate preparation. For a project committed to a specific typeface as part of its visual identity, as this book's own conventions are (TeX Gyre Pagella throughout), this is not a marginal convenience; it is the difference between the font being usable at all and requiring a preparation step that may not succeed.

\hypertarget{fontspec-and-opentype-features}{%
\section{\texorpdfstring{\pkg{fontspec} and OpenType Features}{fontspec and OpenType Features}}\label{fontspec-and-opentype-features}}

\pkg{fontspec} is the interface both LuaLaTeX and XeLaTeX use for font loading, and its feature syntax is worth learning in some depth, since most of what makes a document look professionally typeset rather than merely correctly typeset lives here. \texttt{\textbackslash{}setmainfont\{TeX\ Gyre\ Pagella\}} loads a font by its family name; the same command accepts a bracketed list of OpenType features that changes how that font behaves well beyond simply which glyphs it contains.

Ligatures beyond the basic \texttt{fi} and \texttt{fl} pairs, historical and discretionary ligatures a font may include, are enabled or disabled through the \texttt{Ligatures} feature key. Small capitals, when a font provides a true small-caps design rather than a mechanically shrunk version of the regular capitals, are accessed through \texttt{Letters=SmallCaps} or the \texttt{\textbackslash{}textsc} command once the feature is enabled, and the difference between a genuine small-caps design and a shrunk substitute is visible immediately to anyone who has seen both, which is part of why choosing a font family with real small capitals matters for a document that uses them at all. Old-style figures, numerals that vary in height and descend below the baseline rather than sitting uniformly at cap height, are enabled through \texttt{Numbers=OldStyle}, and suit running prose far better than the lining figures most fonts default to, which were designed to align with tables rather than to sit comfortably inside a sentence.

None of these features change what characters a document contains; they change which glyphs the font substitutes for those characters, entirely at the typesetting layer, which is precisely why \pkg{fontspec}'s feature syntax is worth treating as part of a document's typographic design rather than as an optional flourish. A title page, front matter, and running text that consistently use a font's intended ligatures, small capitals, and old-style figures read as considered and professionally set; the same document with all of a font's default settings left untouched reads as merely competent.

\hypertarget{multiple-font-families-and-mathematics}{%
\section{Multiple Font Families and Mathematics}\label{multiple-font-families-and-mathematics}}

A typeset book rarely uses a single font throughout. \pkg{fontspec} supports loading a separate family for monospaced text (\texttt{\textbackslash{}setmonofont}, typically used for code listings) and a separate family for sans-serif text (\texttt{\textbackslash{}setsansfont}, used for headings in some design conventions or for figures and diagrams), each with its own independent feature set, so that a monospace font chosen for code legibility does not need to share the body text's ligature and old-style-figure settings, which would be actively wrong for code.

Mathematics is the one area where \pkg{fontspec} alone is not sufficient, since math typesetting needs a font specifically designed with the full set of mathematical glyphs, spacing, and sizing variants LaTeX's math mode expects, not merely a text font pressed into service. \pkg{unicode-math}, layered on top of \pkg{fontspec}, provides this, loading a dedicated OpenType math font through \texttt{\textbackslash{}setmathfont} and giving direct Unicode input for mathematical symbols in place of the older, more limited \texttt{amsfonts}/\texttt{amssymb} symbol-by-command approach. A document committed to a particular text typeface throughout, as this book is to TeX Gyre Pagella, generally wants a math font from the same family where one exists, so that text and mathematics read as part of one coherent design rather than as two different typefaces awkwardly sharing a page; where no matching math font exists for a chosen text face, choosing a math font that harmonizes reasonably well in weight and proportion is the next-best option.

\hypertarget{keeping-a-document-engine-agnostic}{%
\section{Keeping a Document Engine-Agnostic}\label{keeping-a-document-engine-agnostic}}

LuaLaTeX and XeLaTeX share \pkg{fontspec} as a common interface, which lets most font-related document code run unchanged under either engine, but the two are not identical underneath that shared interface, and a superuser working across both, or handing a document to a collaborator whose installation defaults to the other engine, needs to know where the differences actually surface.

LuaLaTeX is built on LuaTeX, which exposes a full embedded Lua scripting layer, letting a document call arbitrary Lua code during typesetting; XeLaTeX has no equivalent, so any macro relying on \texttt{\textbackslash{}directlua} or a Lua-based package is LuaLaTeX-only by construction, not merely by convention. XeLaTeX, in turn, has historically had somewhat different and in some cases more mature support for certain complex-script shaping engines through its use of the system's native font-rendering libraries, though LuaTeX's own shaping support has closed most of this gap in recent releases. Microtype's font expansion and protrusion features, which subtly adjust glyph widths to improve justification, behave slightly differently between the two engines because they operate through different underlying mechanisms, which can produce visibly different line breaks for the same source under the two engines even with identical \pkg{fontspec} settings.

For a document meant to remain buildable under either engine, the practical discipline is straightforward: avoid \texttt{\textbackslash{}directlua} and other LuaTeX-specific primitives in document-level code, even when working primarily in LuaLaTeX, unless engine-specific behavior is actually required and is documented as such; test a build under both engines at least once before considering a template or class file finished, rather than assuming shared \pkg{fontspec} syntax guarantees shared output; and treat any package documentation that states an engine requirement as authoritative rather than assuming a feature that works under one engine will silently work under the other. A project that follows this discipline can genuinely move between LuaLaTeX and XeLaTeX as needed; one that does not will eventually break in a way that is difficult to diagnose, because the failure will look like a font problem when its actual cause is an engine-specific primitive buried several packages deep in the document's dependency chain.

The next chapter goes further into what this font infrastructure makes possible beyond ordinary Latin-script typesetting: multilingual documents, custom notation systems, and the construction of entirely new symbol sets built on top of the same OpenType foundation covered here.
